% =============================================================================
%  SPEECH NOTES – Models Section (Slides 5–8)
%  Industrial Component Anomaly Detection – Project Presentation
%  Purpose: Precise speaker guide; what to say, in order, per slide.
% =============================================================================
\documentclass[12pt, a4paper]{article}

\usepackage[utf8]{inputenc}
\usepackage[T1]{fontenc}
\usepackage[margin=2.0cm, top=1.5cm, bottom=2.5cm]{geometry}
\usepackage{lmodern}
\usepackage{microtype}
\usepackage{xcolor}
\usepackage{titlesec}
\usepackage{enumitem}
\usepackage{mdframed}
\usepackage{parskip}
\usepackage{amsmath}
\usepackage{amssymb}
\usepackage{hyperref}
\usepackage{fancyhdr}
\usepackage{booktabs}

% --- Colour palette -----------------------------------------------------------
\definecolor{teal}{RGB}{0,128,128}
\definecolor{accent}{RGB}{180,60,60}
\definecolor{muted}{RGB}{110,110,110}
\definecolor{notegreen}{RGB}{34,139,34}
\definecolor{warnorange}{RGB}{200,100,0}

% --- Section styling ----------------------------------------------------------
\titleformat{\section}[block]
  {\large\bfseries\color{teal}}{\thesection.}{0.5em}{}[\vspace{-0.3em}\rule{\linewidth}{1.2pt}]
\titleformat{\subsection}[block]
  {\normalsize\bfseries\color{accent}}{\thesubsection}{0.5em}{}
\titleformat{\subsubsection}[runin]
  {\small\bfseries\color{muted}}{}{0em}{}[\ ---\ ]

% --- Box styles ---------------------------------------------------------------
\mdfdefinestyle{saybox}{
  backgroundcolor=teal!5, linecolor=teal, linewidth=1.2pt,
  innerleftmargin=8pt, innerrightmargin=8pt,
  innertopmargin=6pt, innerbottommargin=6pt,
  skipabove=6pt, skipbelow=4pt,
}
\mdfdefinestyle{warnbox}{
  backgroundcolor=warnorange!8, linecolor=warnorange, linewidth=1.2pt,
  innerleftmargin=8pt, innerrightmargin=8pt,
  innertopmargin=5pt, innerbottommargin=5pt,
  skipabove=6pt, skipbelow=4pt,
}
\mdfdefinestyle{notebox}{
  backgroundcolor=notegreen!5, linecolor=notegreen, linewidth=1pt,
  innerleftmargin=8pt, innerrightmargin=8pt,
  innertopmargin=5pt, innerbottommargin=5pt,
  skipabove=4pt, skipbelow=4pt,
}

% --- Header / Footer ----------------------------------------------------------
\pagestyle{fancy}
\fancyhf{}
\fancyhead[L]{\small\color{muted}\textbf{Speech Notes} -- Models Section}
\fancyhead[R]{\small\color{muted}Industrial Component Anomaly Detection}
\fancyfoot[C]{\small\thepage}
\renewcommand{\headrulewidth}{0.4pt}

% --- Custom commands ----------------------------------------------------------
\newcommand{\say}[1]{%
  \begin{mdframed}[style=saybox]
  \itshape ``#1''
  \end{mdframed}
}
\newcommand{\warn}[1]{%
  \begin{mdframed}[style=warnbox]
  \textbf{\color{warnorange}$\blacktriangleright$ Watch out:}\ #1
  \end{mdframed}
}
\newcommand{\aside}[1]{%
  \begin{mdframed}[style=notebox]
  \textbf{\color{notegreen}Note to self:}\ #1
  \end{mdframed}
}
\newcommand{\timing}[1]{\hfill{\small\color{muted}[\,\textbf{#1}\,]}}
\newcommand{\pointer}[1]{\noindent\textbf{\color{teal}$\rightarrow$ Point at:} \textit{#1}\\[2pt]}
\newcommand{\transition}[1]{\vspace{4pt}\noindent\textbf{\color{accent}Transition:} \textit{``#1''}\vspace{4pt}}

% =============================================================================
\begin{document}

\begin{center}
  {\Large\bfseries Speech Notes: Models Section}\\[4pt]
  {\normalsize\color{muted} Slides 5--8 $\cdot$ Target: $\approx$6 min total $\cdot$ $\sim$1.5 min/slide average}\\[6pt]
  \begin{tabular}{@{}llr@{}}
    \toprule
    Slide & Title & Target \\
    \midrule
    5 & Model Selection \& Architecture Overview & \textbf{$\approx$2 min} \\
    6 & PatchCore: Transfer Learning \& Architecture & \textbf{$\approx$1.5 min} \\
    7 & Keras CAE -- Autoencoder Learning Process & \textbf{$\approx$1.5 min} \\
    8 & Keras CAE -- End-to-End Pipeline & \textbf{$\approx$1 min} \\
    \bottomrule
  \end{tabular}
\end{center}
\rule{\linewidth}{1.5pt}

% =============================================================================
\section{Slide 5 --- Model Selection \& Architecture Overview \timing{$\approx$2 min}}
% =============================================================================

\subsection{Opening: Entering the Models section}

\say{Now let's talk about the models. Given the task -- unsupervised anomaly detection with zero defect data -- we deliberately chose two architecturally very different approaches, and I want to walk you through why.}

\aside{You just came off the Metrics slide. The audience has AUPIMO in mind. These two models are what AUPIMO will later score.}

% --- LEFT: CAE ---------------------------------------------------------------
\subsection{Left column: Convolutional Autoencoder}

\pointer{the CAE diagram: $x \to$ Encoder $\to z \to$ Decoder $\to \hat{x} \to \Delta$}

\say{The first model is a Convolutional Autoencoder, built completely from scratch. The idea is simple: you take an input image $x$, compress it through a narrow bottleneck $z$, and reconstruct it back as $\hat{x}$. We train exclusively on defect-free images. The model learns to encode what normal looks like. When a defective image arrives at inference, the anomaly gets projected onto the nearest normal point the decoder knows -- and it reconstructs a normal version. The mismatch between input and reconstruction is our anomaly signal.}

\pointer{the score formula below the diagram}

\say{We aggregate that mismatch via Top-K pooling over the per-pixel error map:
\[
  s = \mathrm{TopK}_{0.2\%}\!\bigl(0.84\,(1-\mathrm{SSIM}) + 0.16\,\mathrm{MAE}\bigr)
\]
Top-K at 0.2\% means we take the average of the worst 0.2\% of pixels -- those are your defect candidates. I will come back to the SSIM and MAE choice when we get to the training slide.}

\pointer{the three bullet points: complete control, intuition, trade-offs}

\say{We built this from scratch precisely for the control it gives us: full freedom over the bottleneck dimensions, the loss, the augmentation strategy. The trade-off is a real one -- autoencoders are prone to blurriness or, worse, to learning identity shortcuts if we are not careful.}

% --- RIGHT: PatchCore -------------------------------------------------------
\subsection{Right column: PatchCore + ResNet-18}

\pointer{the PatchCore diagram: $x \to$ ResNet-18 (frozen) $\to \phi(x) \to \mathcal{M} \to d^*$}

\say{The second model is PatchCore, from Intel's Anomalib library. Completely different paradigm: no gradient updates, no training in the traditional sense. Instead, you pass every normal training image through a frozen ResNet-18 backbone, extract mid-level patch feature vectors -- that is $\phi(x)$ -- and store a compressed subset in a memory bank $\mathcal{M}$.}

\say{At inference, the anomaly score for a test patch is just the nearest-neighbour L2 distance to the closest stored normal prototype:
\[
  s(p^*) = \min_{m \in \mathcal{M}} \|\phi(p^*) - m\|_2
\]
If a patch has no close match in memory, it is anomalous. This is closer to database indexing than machine learning.}

\pointer{coreset bullet: ``Compresses the memory bank by 90\%+''}

\say{Anomalib handles coreset subsampling via greedy farthest-point sampling -- the 10\% of vectors that cover the whole feature space best. Our contribution was integrating this into our evaluation protocol and cross-validating it against the CAE.}

\subsection{Why two such different paradigms?}

\say{These two models fail in completely different ways -- and that is precisely why we chose both. A CAE can fail silently on unseen texture distributions. PatchCore fails when memory bank coverage has gaps. Running both gives us orthogonal error signals, and PatchCore also acts as a natural upper-bound baseline: if frozen ImageNet features cannot detect a defect, no learned model will.}

\transition{Let me go into PatchCore first, since it was actually our starting point.}

\vspace{1em}

% =============================================================================
\section{Slide 6 --- PatchCore: Transfer Learning \& Architecture \timing{$\approx$1.5 min}}
% =============================================================================

\subsection{Opening: Why PatchCore was first}

\say{PatchCore was our first model, and the reason is straightforward: it was the most intuitive entry point into anomaly detection without any defect data. Anomalib gave us a fully working pipeline out of the box -- we could evaluate it on day one, before writing a single line of custom model code. That made it an invaluable baseline and an anchor point for everything that followed.}

% --- Left column: Flowchart -------------------------------------------------
\subsection{Left column: Transfer Learning flowchart}

\pointer{the four-step vertical flowchart}

\say{Walk through top to bottom.

\textbf{Normal Factory Images:} zero defect labels needed, Day-1 deployable.

\textbf{Frozen Feature Extractor:} ResNet-18, pre-trained on ImageNet, weights fully locked. We use layers 2 and 3 specifically -- not the final classification layer which encodes \textit{cat vs truck} semantics, but the mid-level layers that encode edges, textures, and local structures. Those transfer perfectly to surface defect domains.

\textbf{Coreset Memory Bank:} greedy farthest-point sampling compresses all patch vectors down to $\sim$10\%, keeping maximum coverage of the feature space manifold.

\textbf{Inference:} nearest-neighbour distance to $\mathcal{M}$. Large distance means anomalous.}

% --- Right column: Results + bullets ----------------------------------------
\subsection{Right column: Results and key properties}

\pointer{the four heatmap images: wood, zipper, transistor, toothbrush}

\say{These are PatchCore's heatmaps across four MVTec categories. White overlay is the ground truth; green is the predicted mask at the 99th percentile threshold. The match is strong even for visually challenging defects like the liquid contamination on wood.}

\pointer{Key Features and ``Why it excels'' bullet lists}

\say{Two properties I want to highlight. First, cold-start eliminated -- no defect data collection before deployment. Second, IP protection: the memory bank stores embeddings, not raw images -- it is not reversible.

On why it excels: the frozen backbone is critical. If you fine-tuned ResNet on normal images only, you would collapse the embedding space to the factory distribution and lose exactly the generalisation that makes PatchCore sensitive to subtle defects.}

\transition{Now let us look at our custom model -- the Keras CAE -- and how we built it up from scratch, failure mode by failure mode.}

\vspace{1em}

% =============================================================================
\section{Slide 7 --- Keras CAE: Autoencoder Learning Process \timing{$\approx$1.5 min}}
% =============================================================================

\subsection{Opening: A story of incremental decisions}

\say{This slide tells the development story of the CAE. We did not arrive at the final architecture on the first try. Every design decision here was motivated by a concrete problem we ran into.}

% --- Left column: Concept ---------------------------------------------------
\subsection{Left column: The core concept}

\pointer{Training Phase / Inference Phase bullets}

\say{The concept is simple: train on normals, fail to reconstruct defects at inference. But the \textit{how} is what matters -- and that is the right column.}

% --- Right column: Design evolution -----------------------------------------
\subsection{Right column: Key design decisions}

\subsubsection{Activation: ReLU $\to$ (LeakyReLU) $\to$ ELU}

\say{We started with ReLU. The dying neuron problem hit us fast: neurons in the bottleneck locked to hard zero and stopped contributing gradients entirely. On our relatively limited dataset -- only 200 to 300 normal images per category -- we cannot afford to lose representation capacity. With ReLU, any negative input is hard zeroed out. With ELU, the chance of hitting exact zero is virtually zero: for negative inputs, it smoothly follows $\alpha(e^x - 1)$, ensuring continuous, non-zero gradient flow so every parameter keeps learning even across small batches. We also tested Leaky ReLU, but its sharp kink at zero introduced noise into fine-grained reconstruction gradients. ELU is continuously differentiable, saturates gracefully for large negative inputs, and naturally pushes mean activations toward zero. That made it the optimal choice for our small-data regime.}

\subsubsection{BatchNorm and bias}

\say{Early on we used bias terms and no BatchNorm. Once we added BatchNormalisation, we removed the bias. This is by design: BatchNorm already shifts activations through its own learnable parameter $\beta$. A separate bias is redundant -- it is absorbed into $\beta$ -- so \texttt{use\_bias=False} keeps the model clean. BatchNorm also stabilises training by reducing internal covariate shift, which on small datasets like MVTec is very valuable.}

\subsubsection{Optimizer: Adam $\to$ AdamW}

\say{Standard Adam couples weight decay with the adaptive learning rate, which under-regularises on small datasets. AdamW decouples them: the gradient step and the weight decay apply independently. With only a few hundred normal training images per category, that stronger regularisation matters.}

\subsubsection{MIM: preventing the identity shortcut}

\pointer{``Masked Image Modelling: 25\% of patches randomly masked'' bullet}

\say{Finally, Masked Image Modelling: we randomly zero 25\% of input patches before the encoder sees them. Without this, the model can learn a trivial identity map and reconstruct anomalies just as well as normals -- total failure. MIM forces it to infer missing content from context and build a genuine structural model of what normal looks like.}

\subsubsection{Training loss vs. inference score -- Literature vs. Implementation}

\say{Here is a technical nuance worth clarifying regarding literature versus implementation. In the seminal paper by Bergmann et al.\ (2019, VISAPP), they evaluated SSIM with both $L_1$ and $L_2$, and specifically recommended $\text{SSIM} + L_1$ (MAE) for both training and evaluation.

In our project:
\begin{itemize}[topsep=2pt,itemsep=2pt]
  \item During \textbf{training}, we implemented $\text{SSIM} + \text{MSE}$ ($L_2$):
  \[
    \mathcal{L}_\text{train} = 0.84\,(1-\mathrm{SSIM}) + 0.16\,\mathrm{MSE}
  \]
  MSE provided smooth, stable gradient convergence in our Keras pipeline.
  \item At \textbf{inference}, our pixel anomaly map strictly aligns with Bergmann's formula, using $\text{SSIM} + \text{MAE}$ ($L_1$):
  \[
    \text{Score} = \mathrm{TopK}_{0.2\%}\!\bigl(0.84\,(1-\mathrm{SSIM}) + 0.16\,\mathrm{MAE}\bigr)
  \]
  Why use MAE for inference? Because squaring errors with $L_2$ at test time penalises large errors quadratically, blowing out bright spots and distorting localization heatmaps. Linear $L_1$ error preserves the true proportional contrast of defects.
\end{itemize}}

\warn{If an examiner asks: \textit{``Didn't Bergmann recommend SSIM + L1 for both?''} \\
\textbf{Answer:} \textit{``Yes, absolutely! Bergmann et al.\ (2019) found SSIM + $L_1$ best overall. In our Keras training implementation, we coupled SSIM with MSE ($L_2$) because it offered smooth gradient convergence during training, but for inference and error map generation, we stayed true to Bergmann's linear $L_1$ formulation so the pixel error map isn't distorted quadratically by squared errors.''}}

\transition{Let me now show you how all these pieces come together in the full automated pipeline.}

\vspace{1em}

% =============================================================================
\section{Slide 8 --- Keras CAE: End-to-End Pipeline \timing{$\approx$1 min}}
% =============================================================================

\subsection{Step-by-step walkthrough}

\pointer{the 9-step TikZ pipeline diagram -- three rows}

\say{The pipeline runs left to right across three rows. Let me go through it quickly.}

\vspace{4pt}
\noindent\textbf{Row 1 -- Data Preparation (Steps 1--3):}

\say{Raw MVTec data goes through preprocessing -- Otsu plus Canny edge detection for foreground masking, CLAHE contrast enhancement, and blur -- then branches on category type. Textures get rotation and flip augmentations; objects get colour jitter and noise. This category-aware branching matters because texture datasets are rotationally invariant while object datasets are not.}

\vspace{4pt}
\noindent\textbf{Row 2 -- Model Building and Training (Steps 4--6):}

\say{Step 4 applies Masked Image Modelling -- 25\% of patches zeroed. Step 5 builds the CAE with ELU, BatchNorm, AdamW. Step 6 trains with the \textbf{SSIM+MSE loss} -- this is the training loss, not the anomaly score.}

\vspace{4pt}
\noindent\textbf{Row 3 -- Inference and Evaluation (Steps 7--9):}

\say{Step 7 scores images: frozen weights, per-pixel \textbf{SSIM+MAE error map}, Top-K pooling to a single image-level score. Step 8 applies the adaptive threshold -- calibrated on the 15\% normal validation partition using quantile or Mahalanobis distance. This gives a Day-1 deployable false-alarm-rate guarantee with zero defect knowledge. Step 9 evaluates on the locked test set: AUROC and AUPIMO.}

\pointer{the Optuna caption at the bottom}

\say{The whole pipeline is orchestrated by Optuna -- Bayesian optimisation over augmentation strategy, latent dimensions, crop size -- targeting Pixel AUPIMO as the objective. We target AUPIMO rather than AUROC because AUPIMO evaluates only the operationally relevant region below false-alarm rate $10^{-4}$. Optimising AUROC could improve the broad curve while degrading exactly the narrow region that matters in a real factory.}

\subsection{Closing the Models section}

\say{So to summarise the Models section: PatchCore is our high-accessibility no-training baseline that operationalises frozen ImageNet features. The Keras CAE is our from-scratch model where every decision -- ELU, BatchNorm, MIM, AdamW, and the loss-versus-score distinction -- was driven by a concrete failure mode we needed to address. In the next section I will show you how we evaluated both under the same zero-leakage protocol.}

\vspace{2em}
\rule{\linewidth}{0.8pt}
{\small\color{muted}\textit{End of Models section speech notes.}}

\end{document}
